%--------------------------------------------------------------------------------
%
%
%--------------------------------------------------------------------------------
\pagebreak
\section*{FDCA — Flush Data Cache}
\addcontentsline{toc}{section}{FDCA — Flush Data Cache}

%--------------------------------------------------------------------------------
\begin{iPageDescEntry}{Syntax}
\texttt{fdca <vAdr>} \\
\end{iPageDescEntry}

%--------------------------------------------------------------------------------
\begin{iPageDescEntry}{Description}

The \texttt{FDCA} command flushes a single cache line from the data cache.
The cache line is selected using the specified virtual address.

For processors with a unified instruction and data cache, the command refers to the unified cache.

This command requires the current window to be a data cache window.

\end{iPageDescEntry}

%--------------------------------------------------------------------------------
\begin{iPageDescEntry}{Example}

The following command flushes the data cache line corresponding to the virtual
address \texttt{0\_0001\_ffff\_0000}.

\begin{lstlisting}[style=iPageOpStyle]
(10) -> fdca 0_0001_ffff_0000
(11) ->
\end{lstlisting}

\end{iPageDescEntry}

%--------------------------------------------------------------------------------
\begin{iPageDescEntry}{Notes}
None.
\end{iPageDescEntry}
